\newpage
\thispagestyle{empty}
\mbox{}
\newpage

%%%%%%%%%%%%%%%%%%%%%%%%%%%%%%%%%%%%%%%%%%%%%%%%%%%%%%%%%%%%

\chapter{Contributions to the scientific community}
\chaptermark{Contributions to the scientific community}
\label{chap:contrib_commu}

This side chapter is intended to list, and shortly describe, the various contributions of my thesis and side thesis activities to the scientific community.
Here are presented the corpus of publications for which I contributed.

\section*{Peer-reviewed publications}
    \subsection*{Thesis research outputs}
    \begin{itemize}
        \item Published in \textit{Atmospheric Chemistry and Physics}:
            \textbf{Da Silva, A.}, Marelle, L., Raut, J.-C., Gramlich, Y., Siegel, K., Haslett, S. L., Mohr, C., and
Thomas, J. L. (2025): How to trace the origins of short-lived atmospheric species: an Arctic
            example, Atmospheric Chemistry and Physics. DOI: \url{https://doi.org/10.5194/acp-25-5331-2025}

        \item blablablah

    \end{itemize}
    \subsection*{Scientific collaborations}
    \begin{itemize}
        \item Published in \textit{Atmospheric Chemistry and Physics}:
        Lapere, R., Thomas, J. L., Favier, V., Angot, H., Asplund, J., Ekman, A. M. L., Marelle, L, Raut, J-C, \textbf{Da Silva, A}, Wille, J. D., Zieger, P. (2024). Polar aerosol atmospheric rivers: Detection, characteristics, and potential applications. Journal of  Geophysical Research: Atmospheres, 129, e2023JD039606. \url{https://doi.org/10.1029/2023JD039606}.

        \item blablablah 

    \end{itemize}

\section*{Conferences, colloquia, and workshops}
    \subsection*{2023}
    \begin{itemize}
        \item EGU general assembly 2023, Vienna, oral presentation:

        \textbf{Da Silva, A.}, Marelle, L., and Raut, J.-C.: Seeking the origins of Arctic ice nucleating particles with FLEXPART-WRF, EGU General Assembly 2023, Vienna, Austria, 24–28 Apr 2023, EGU23-6392, \url{https://doi.org/10.5194/egusphere-egu23-6392, 2023}. 

        \item blablablah 
    \end{itemize}


\section*{International and national science projects}
    \begin{itemize}
        \item Climate Relevant interactions and feedbacks: the key role of sea ice and Snow in the polar and global climate system (CRiceS).

            The CRiceS project focuses on improving model predictions of the role of polar processes in the climate system that consists of the oceans, ice and snow cover, and the atmosphere. It is crucial to understand the role of the polar processes, such as feedback loops, in polar and global climate. One of the main ways scientists can improve our understanding of environmental change is to combine knowledge from different disciplines in a coordinated way.

            My work contributed to the work package 2 (Improved ocean-ice/snow-atmosphere processes in models) and to the core theme 2 (Aerosols and Clouds).

        \item  blablabla
    \end{itemize}


\section*{Teaching activities}
    \subsection*{Sorbonne Université}
    I took part in six different courses, teaching students from first year to master level, within two departments of the university (\textit{Earth's sciences} and \textit{Physics}). 
    Below stands the detailed list of those courses. 
    
    \begin{itemize}
        \item[] \textbf{Undergraduate}
        \begin{itemize}
            \item \textit{Terre, Climat, Environnement}: tutorial leading and elaboration of presentations topics for the new maquette of the course.
        \end{itemize}
        %\item[] \textbf{L2} 
        \begin{itemize}
            \item \textit{Le changement climatique et ses enjeux}: supervising of students talks about climate-society relations.
        \end{itemize}
        %\item[] \textbf{L3}  
        \begin{itemize}
            \item \textit{Matlab modelling project}: supervising of modelling projects for five groups of students in geology.
            \item \textit{Physique des millieux continus} : charge des TDs et des résolutions de problème (RP).
        \end{itemize}
        \item[] \textbf{Graduate}
        \begin{itemize}
            \item \textit{Statistics} : tutorial leading and project supervising in advanced statistics with R.
            \item \textit{Rayonnement et télédetection} : supervising of the lidar practical work session.
        \end{itemize}
    \end{itemize}

    \subsection*{Ecole nationale supérieure des études avancées (ENSTA) Paris}
    I have been tutorial leader for the transdisciplinary course named \textit{Analyse systémique de l'Anthropocène} (Systemic analysis of the Anthropocene), for ingineering school students.

    \subsection*{European winter school ERCA 2025}
    I have been project group leader for six international PhD students. I supervised the group that focused on aerosol-cloud interactions. 
    The objective I set up was to highlight the effect of ice nucleating particles on Arctic clouds thanks to principal component analysis of in situ measurements, and to empirical orthogonal functions computed on reanalysis data (ERA5 and CAMS).

\section*{Climactions IPSL}
    Under the name of \textit{Climactions}, gather researchers of the Institut Pierre-Simon Laplace (IPSL) laboratories in a dynamic of thinking how to reduce the environmental footprint of their research. Since its launch in 2019, Climactions has expanded its objectives and shifted its focus: it is no longer only concerned with the environmental impact of research, but also with how to make research practices resilient in a society increasingly affected by the consequences of \tempc{2} global warming.
    I have participated to this dynamic through different activities presented below.

    \subsection*{Working group research core activities}
    During the first two years of my PhD, I hosted the IPSL working group about the evolution of the observational activities of the climate sciences community.
    The group gathered researchers of various fields among the IPSL community. Our monthly meeting aimed to communicate and think about the levers to make scientific observation activities more resilient and sufficient: from the conception of the next generation of the French fleet of oceanographic ships, to the future of space observation, by way of Antarctic polar bases. 
    The outline of our work can be found at \href{https://climactions.ipsl.fr/groupes-de-travail/gt2-observations/}{Climactions' working group on observations activities}\footnote{\url{https://climactions.ipsl.fr/groupes-de-travail/gt2-observations/}}.


    \subsection*{Villarceaux Agreement and tool box}
    I actively participated to the elaboration of a panel of almost \num{50} measures (named the tool box), intended to help the IPSL laboratories reducing their environmental footprint. They also aimed to help the laboratories getting more resilient against budget cuts and various kinds of physical restrictions that public research may face in the near future. 

    In parallel, we set up the \textit{Villarceaux Agreement} which\---- in the manner of the Paris Agreement\----states the recognition of the need for our institute to change, reduce and adapt its activities, as well as the commitment to do so in order to preserve good research conditions for the present and next generations of scientists.
    Today, the Agreement has been signed by seven laboratories of the IPSL. The text can be found here: \href{https://cloud.ipsl.fr/index.php/s/e7YaN7JEqw5r9nK?openfile=true}{Villarceaux agreement text}\footnote{\url{https://cloud.ipsl.fr/index.php/s/e7YaN7JEqw5r9nK?openfile=true}}.

    \subsection*{Climactions at LATMOS}
    I have participated to the local LATMOS Climactions group. With Lola Falletti (head of the LATMOS Climactions group), we organised a special general assembly to present four measures that have been then voted by all the laboratories members (\href{https://climactions.ipsl.fr/groupes-de-travail/latmos-climactions/}{LATMOS Climactions web page}\footnote{\url{https://climactions.ipsl.fr/groupes-de-travail/latmos-climactions/}}).
    The next year, I organised a webinar to prepare the IPSL Transform'Actions week, when all IPSL laboratories gathered to discuss and choose measures from the tool box that best suit their activities.
    During that week, I led the LATMOS discussions.
    All along the three years of my PhD, I have made the link between the LATMOS directors and the IPSL head about Climactions affairs.


\section*{Science communication and awarness raising}
    \subsection*{\textit{Ma thèse en 180 secondes} 2023 contest}
    I participated to the francophone contest \textit{Ma thèse en 180 secondes}, where PhD candidates from every fields have 180 seconds to present their thesis to a non-specialist audience. I did the Sorbonne Université final. My appearance has been subtitled \url{https://www.youtube.com/watch?v=pGzU3e5XTz0}{here}\footnote{\url{https://www.youtube.com/watch?v=pGzU3e5XTz0}}.

    \subsection*{High school teaching}
        \begin{itemize}
            \item Interventions about global warming and climate sciences in high school classes (\textit{2\up{nde}}, \textit{1\up{ère}} and \textit{terminales}).
            \item Atmospheric physics workshop with \textit{terminales} students specialized in mathematics. 
        \end{itemize}


    \subsection*{RADICAL}
    On the impulsion of François Ravetta, former director of LATMOS, his PhD student Thomas Lesigne and have held a scientific mediation mission about the RADiomètre Infrarouge pour la Course Au Large (RADICAL), a small infrared radiometer that the LATMOS installed on a racing sail boat. The project is resumed in the Radical \href{https://drive.google.com/file/d/1CV0wxgG6JPSrtKX6_IJspliK54sUBc2E/view?usp=sharing}{booklet}\footnote{\url{https://drive.google.com/file/d/1CV0wxgG6JPSrtKX6_IJspliK54sUBc2E/view?usp=sharing}} and \href{https://drive.google.com/file/d/17i1LEQ7NheX6M8gCM3Y0SALK7NreZvb9/view?usp=sharing}{kakemono poster}\footnote{\url{https://drive.google.com/file/d/17i1LEQ7NheX6M8gCM3Y0SALK7NreZvb9/view?usp=sharing}} we produced for the occasion. 


%\section*{Lab life}
%...





